\documentclass{article}
\usepackage{amsmath,amssymb}
\usepackage{xurl}
\usepackage[colorlinks=true,linkcolor=blue,citecolor=blue,urlcolor=blue]{hyperref}
% Shared commands for notes in docs/paper-gaps/.

\DeclareUnicodeCharacter{2080}{\ifmmode{}_{0}\else\textsubscript{0}\fi}
\DeclareUnicodeCharacter{2081}{\ifmmode{}_{1}\else\textsubscript{1}\fi}
\DeclareUnicodeCharacter{2082}{\ifmmode{}_{2}\else\textsubscript{2}\fi}
\DeclareUnicodeCharacter{2099}{\ifmmode{}_{n}\else\textsubscript{n}\fi}

\newcommand{\Lean}{\textsc{Lean}}
\ExplSyntaxOn
\NewDocumentCommand{\leanid}{m}{
  \ifmmode
    \textsf{\detokenize{#1}}
  \else
    \group_begin:
    \urlstyle{sf}
    \tl_set:Nn \l_tmpa_tl {#1}
    \tl_replace_all:Nnn \l_tmpa_tl {\_} {_}
    \exp_args:NV \path \l_tmpa_tl
    \group_end:
  \fi
}
\ExplSyntaxOff
\providecommand{\tr}{\operatorname{tr}}
\newcommand{\tnleanIssueBase}{https://github.com/LionSR/TNLean/issues/}
\newcommand{\tnleanPrBase}{https://github.com/LionSR/TNLean/pull/}
\newcommand{\tnleanDocBase}{https://sirui-lu.com/TNLean/docs/}
\newcommand{\ghissue}[1]{\href{\tnleanIssueBase#1}{GitHub issue \##1}}
\newcommand{\ghpr}[1]{\href{\tnleanPrBase#1}{PR \##1}}
\newcommand{\doclink}[1]{\href{\tnleanDocBase#1.html}{\path{#1}}}

% Dirac notation and the identity operator, as in blueprint/src/macros/common.tex.
% \providecommand keeps note-local definitions authoritative.
\usepackage{dsfont}
\providecommand{\ket}[1]{|#1\rangle}
\providecommand{\bra}[1]{\langle #1|}
\providecommand{\braket}[2]{\langle #1 | #2 \rangle}
\providecommand{\ketbra}[2]{|#1\rangle\!\langle #2|}
\providecommand{\Id}{\mathds{1}}
\providecommand{\one}{\mathds{1}}
\providecommand{\C}{\mathbb{C}}
\providecommand{\MN}[1]{M_{#1}(\C)}
\providecommand{\MD}{M_D(\C)}

% Verdict marker: \gapnote{<kind>}{<status>}. Kind is the policy
% classification (clarification, local-correction, scope-restriction,
% unfaithful, false-source, open-gap); status is open, wip (elimination
% underway), resolved, or historical. Machine-read by the site index;
% prints a verdict line.
\newcommand{\gapnote}[2]{\par\noindent\textbf{Verdict:} #1 (#2).\par}

\title{Threshold Truncation in the One-Dimensional Matrix Product State Approximation}
\author{}
\date{October 18, 2026}
\begin{document}
\maketitle
\gapnote{scope-restriction}{resolved}

\section{The source statement}

Proposition~7.1 of the one-dimensional note~\cite{TNLean2026OneDimensional}
(label \texttt{prop:VC}, file \texttt{part2a\_mps\_information.tex},
lines~39--62) is the classical sequential truncation of Verstraete and
Cirac. Let $\Omega$ be a unit vector on the chain $\Lambda=\{1,\ldots,L\}$
and let $D$ be an arbitrary natural number. For $1\le k\le L-1$ let
$\Pi_k$ be the projection of $\mathcal H_{[1,k]}$ onto the span of the
$\min\{D,\dim\mathcal H_{[1,k]}\}$ eigenvectors of $\rho_{[1,k]}$ with the
largest eigenvalues (the identity when $D\ge\dim\mathcal H_{[1,k]}$), let
$\epsilon_k=\langle\Omega,(\mathbb 1-\Pi_k\otimes\mathbb 1)\Omega\rangle$ be
the discarded mass, and put
\begin{align}
    \Phi
     & = (\Pi_1\otimes\mathbb 1)(\Pi_2\otimes\mathbb 1)\cdots
    (\Pi_{L-1}\otimes\mathbb 1)\,\Omega .
    \label{eq:tnlean_one_dimensional_threshold_truncation_truncated_vector}
\end{align}
Then $\Phi$ is a matrix product state with open boundary conditions and bond
dimension at most $D$, and $\|\Phi-\Omega\|\le\sum_k\sqrt{\epsilon_k}$.

The proposition quantifies over the bond dimension $D$: the truncation at
every cut keeps the $D$ largest Schmidt weights, whatever the spectrum, and
the error is expressed through the masses $\epsilon_k=\epsilon_k(D)$ beyond
the $D$-th eigenvalue. Theorem~B is then derived (the two-line route,
label \texttt{thm:Bshort}, lines~80--92) by choosing a slack $w$ from the
tail bound and taking $D=\lceil e^{C_{\mathrm{int}}+w}\rceil$, so that the
eigenvalues at least $e^{-(S(X)+w)}$ are among the $D$ largest and
$\epsilon_k(D)$ is at most the surprisal tail beyond $w$.

\section{The former threshold restriction}

The earlier formalization established only the \emph{threshold form} of the sequential
truncation.\footnote{Declarations:
    \leanid{TNLean.MPS.AreaLaw.sequentialTruncation},
    \leanid{TNLean.MPS.AreaLaw.norm_sub_sequentialTruncation_le},
    \leanid{TNLean.MPS.AreaLaw.rank_cutMatrix_sequentialTruncation_le},
    \leanid{TNLean.MPS.AreaLaw.exists_sequentialTruncation}.} For a unit vector
$\psi$ on $L$ sites and a real slack $w$, the regional projector $\Pi_k$ at
the cut $k$ is the spectral projector of the reduced state of the first $k$
sites onto the eigenvalues at least $e^{-(S_k+w)}$, where $S_k$ is the entropy
of that reduced state; its rank is at most $e^{S_k+w}$, and the mass it
discards, $\|(\mathbb 1-\Pi_k\otimes\mathbb 1)\psi\|^2$, is at most the
two-sided surprisal tail of the reduced state beyond $w$.\footnote{Declarations:
    \leanid{Matrix.PosSemidef.rank_spectralProjectionGE_exp_le},
    \leanid{Matrix.PosSemidef.trace_complement_spectralProjectionGE_le_surprisalTail}.}
The vector $\Phi=(\Pi_1\otimes\mathbb 1)\cdots(\Pi_{L-1}\otimes\mathbb 1)\psi$
has cut rank at most $\max\{1,e^{S_k+w}\}$ across every cut and satisfies
$\|\psi-\Phi\|\le\sum_{k=1}^{L-1}\sqrt{\mathrm{tail}_k(w)}$. Together with the
passage from bounded cut ranks to an open-boundary matrix product
state,\footnote{Declaration:
    \leanid{TNLean.MPS.AreaLaw.isOpenMPS_of_cutMatrix_rank_le}.} this is exactly
what the two-line route to Theorem~B uses, and Theorem~B is proved from
it.\footnote{Declaration:
    \leanid{TNLean.MPS.AreaLaw.exists_isOpenMPS_norm_sub_le}.}

The former restriction concerned the choice of projectors. The source projects onto a
prescribed number $D$ of Schmidt vectors at every cut and quantifies over
$D$; the threshold statement projects onto the eigenvalues above a threshold
fixed by the entropy of the cut and a common slack $w$, and quantifies over
$w$. For a prescribed $D$ the two families of projectors need not agree.
The threshold lemma alone therefore does not give the source proposition
for an arbitrary $D$, or its error bound $\sum_k\sqrt{\epsilon_k(D)}$.
The threshold form is
stronger in one respect, the rank bound $e^{S_k+w}$ is exhibited rather than
assumed, and weaker in the other, the bond dimension is not a free parameter.

\section{Resolution by ordered-span truncation}

The completion treats arbitrary ordered orthonormal bases at the internal
cuts. For each cut $k$, let $N_k=\dim\mathcal H_{[1,k]}$, choose an ordered
orthonormal basis $b_{k,0},\ldots,b_{k,N_k-1}$, and let $\Pi_k$ be the
orthogonal projection onto the span of its first $\min\{D,N_k\}$ vectors.
The ordered product
$\Phi=(\Pi_1\otimes\mathbb 1)\cdots(\Pi_{L-1}\otimes\mathbb 1)\psi$
has an open-boundary matrix product representation with internal bond
dimensions at most $D$, and
\begin{align}
    \|\Phi-\psi\|
     & \le\sum_{k=1}^{L-1}
    \sqrt{\operatorname{Re}\langle\psi,(\mathbb 1-\Pi_k\otimes\mathbb 1)\psi\rangle}.
    \label{eq:tnlean_one_dimensional_threshold_truncation_ordered_span_error}
\end{align}
The statement holds for every vector and every such basis family, for all
natural $L,q,D$.\footnote{Declarations:
    \leanid{TNLean.MPS.AreaLaw.prefixBasisProjection},
    \leanid{TNLean.MPS.AreaLaw.isOpenMPS_sequentialTruncation_prefixBasisProjection}.}
For $L\le1$ the product is empty, so $\Phi=\psi$ and the error sum is zero.

In particular, choose a unit vector and any orthonormal eigenbasis of each
marginal, ordered by nonincreasing eigenvalue. These leading spans are the
source spans, including every permissible choice within a repeated
eigenspace, and the displayed expectations are $\epsilon_k(D)$. This
specialization proves Proposition~7.1 and resolves the former restriction.
The note is retained to distinguish the two truncation statements. The
threshold theorem continues to supply the argument for Theorem~B; its
projectors and the prescribed-dimension projectors need not coincide.
Both estimates concern the projected vector before normalization.

\bibliographystyle{alpha}
\bibliography{references}
\end{document}
